% Options for packages loaded elsewhere
\PassOptionsToPackage{unicode}{hyperref}
\PassOptionsToPackage{hyphens}{url}
%
\documentclass[
]{article}
\usepackage{amsmath,amssymb}
\usepackage{iftex}
\ifPDFTeX
  \usepackage[T1]{fontenc}
  \usepackage[utf8]{inputenc}
  \usepackage{textcomp} % provide euro and other symbols
\else % if luatex or xetex
  \usepackage{unicode-math} % this also loads fontspec
  \defaultfontfeatures{Scale=MatchLowercase}
  \defaultfontfeatures[\rmfamily]{Ligatures=TeX,Scale=1}
\fi
\usepackage{lmodern}
\ifPDFTeX\else
  % xetex/luatex font selection
\fi
% Use upquote if available, for straight quotes in verbatim environments
\IfFileExists{upquote.sty}{\usepackage{upquote}}{}
\IfFileExists{microtype.sty}{% use microtype if available
  \usepackage[]{microtype}
  \UseMicrotypeSet[protrusion]{basicmath} % disable protrusion for tt fonts
}{}
\makeatletter
\@ifundefined{KOMAClassName}{% if non-KOMA class
  \IfFileExists{parskip.sty}{%
    \usepackage{parskip}
  }{% else
    \setlength{\parindent}{0pt}
    \setlength{\parskip}{6pt plus 2pt minus 1pt}}
}{% if KOMA class
  \KOMAoptions{parskip=half}}
\makeatother
\usepackage{xcolor}
\usepackage[margin=1in]{geometry}
\usepackage{graphicx}
\makeatletter
\def\maxwidth{\ifdim\Gin@nat@width>\linewidth\linewidth\else\Gin@nat@width\fi}
\def\maxheight{\ifdim\Gin@nat@height>\textheight\textheight\else\Gin@nat@height\fi}
\makeatother
% Scale images if necessary, so that they will not overflow the page
% margins by default, and it is still possible to overwrite the defaults
% using explicit options in \includegraphics[width, height, ...]{}
\setkeys{Gin}{width=\maxwidth,height=\maxheight,keepaspectratio}
% Set default figure placement to htbp
\makeatletter
\def\fps@figure{htbp}
\makeatother
\usepackage{soul}
\setlength{\emergencystretch}{3em} % prevent overfull lines
\providecommand{\tightlist}{%
  \setlength{\itemsep}{0pt}\setlength{\parskip}{0pt}}
\setcounter{secnumdepth}{5}
\usepackage{booktabs}
\usepackage{booktabs}
\usepackage{longtable}
\usepackage{array}
\usepackage{multirow}
\usepackage{wrapfig}
\usepackage{float}
\usepackage{colortbl}
\usepackage{pdflscape}
\usepackage{tabu}
\usepackage{threeparttable}
\usepackage{threeparttablex}
\usepackage[normalem]{ulem}
\usepackage{makecell}
\usepackage{xcolor}
\ifLuaTeX
  \usepackage{selnolig}  % disable illegal ligatures
\fi
\IfFileExists{bookmark.sty}{\usepackage{bookmark}}{\usepackage{hyperref}}
\IfFileExists{xurl.sty}{\usepackage{xurl}}{} % add URL line breaks if available
\urlstyle{same}
\hypersetup{
  pdftitle={Untitled},
  hidelinks,
  pdfcreator={LaTeX via pandoc}}

\title{Untitled}
\author{}
\date{\vspace{-2.5em}2025-01-30}

\begin{document}
\maketitle

{
\setcounter{tocdepth}{2}
\tableofcontents
}
\hypertarget{situation-du-cours-deau}{%
\section{\texorpdfstring{\ul{Situation du cours
d'eau}}{Situation du cours d'eau}}\label{situation-du-cours-deau}}

\begin{longtable}[l]{>{\raggedleft\arraybackslash}p{2cm}>{\raggedleft\arraybackslash}p{2cm}>{\raggedleft\arraybackslash}p{2cm}>{\raggedleft\arraybackslash}p{2cm}>{\raggedleft\arraybackslash}p{2cm}>{\raggedleft\arraybackslash}p{2cm}>{\raggedleft\arraybackslash}p{2cm}}
\toprule
Code de la Masse d'eau & Nom du cours d'eau & Longueur du cours d'eau (en km) & Altitute du point le plus en aval (en m) & Altitude du point le plus en amout (en m) & Dénivelé naturel (en m) & Pente naturelle (en ‰)\\
\midrule
\endfirsthead
\multicolumn{7}{@{}l}{\textit{(continued)}}\\
\toprule
Code de la Masse d'eau & Nom du cours d'eau & Longueur du cours d'eau (en km) & Altitute du point le plus en aval (en m) & Altitude du point le plus en amout (en m) & Dénivelé naturel (en m) & Pente naturelle (en ‰)\\
\midrule
\endhead

\endfoot
\bottomrule
\endlastfoot
\cellcolor{gray!15}{FRHR163 - FRHR162} & \cellcolor{gray!15}{la béthune} & \cellcolor{gray!15}{64.39} & \cellcolor{gray!15}{3.54} & \cellcolor{gray!15}{203.7} & \cellcolor{gray!15}{200.16} & \cellcolor{gray!15}{3.11}\\*
\end{longtable}

\hypertarget{nombre-douvrages}{%
\section{\texorpdfstring{\ul{Nombre
d'ouvrages}}{Nombre d'ouvrages}}\label{nombre-douvrages}}

\begin{longtable}[l]{>{\raggedleft\arraybackslash}p{4cm}>{\raggedleft\arraybackslash}p{4cm}>{\raggedleft\arraybackslash}p{4cm}>{\raggedleft\arraybackslash}p{4cm}}
\toprule
Nom du cours d'eau & Nombre d'ouvrages & Nombre d'ouvrages sur le cours principal & Densité d'ouvarge sur le cours principal par km\\
\midrule
\endfirsthead
\multicolumn{4}{@{}l}{\textit{(continued)}}\\
\toprule
Nom du cours d'eau & Nombre d'ouvrages & Nombre d'ouvrages sur le cours principal & Densité d'ouvarge sur le cours principal par km\\
\midrule
\endhead

\endfoot
\bottomrule
\endlastfoot
\cellcolor{gray!15}{la béthune} & \cellcolor{gray!15}{34} & \cellcolor{gray!15}{20} & \cellcolor{gray!15}{0.31}\\*
\end{longtable}

\hypertarget{ruxe9partition-des-hauteurs-de-chute}{%
\section{\texorpdfstring{\ul{Répartition des hauteurs de
chute}}{Répartition des hauteurs de chute}}\label{ruxe9partition-des-hauteurs-de-chute}}

\includegraphics{test_files/figure-latex/Barplot classe_hauteur et etat-1.pdf}

Ce graphique présente le nombre d'ouvrage par intervals de hauteur de
chute présent sur le cours d'eau. Nous pouvons également observer la
proportion d'ouvrage en fonction de leur état.

\hypertarget{uxe9tat-des-ouvrage-sur-le-cours-principal}{%
\section{\texorpdfstring{\ul{État des ouvrage sur le cours
principal}}{État des ouvrage sur le cours principal}}\label{uxe9tat-des-ouvrage-sur-le-cours-principal}}

\begin{longtable}[l]{rrr}
\toprule
État & Nombre d'ouvrages & Proportion (en \%)\\
\midrule
\endfirsthead
\multicolumn{3}{@{}l}{\textit{(continued)}}\\
\toprule
État & Nombre d'ouvrages & Proportion (en \%)\\
\midrule
\endhead

\endfoot
\bottomrule
\endlastfoot
\cellcolor{gray!15}{Détruit entièrement} & \cellcolor{gray!15}{4} & \cellcolor{gray!15}{0.20}\\
Détruit partiellement & 5 & 0.25\\
\cellcolor{gray!15}{Existant} & \cellcolor{gray!15}{11} & \cellcolor{gray!15}{0.55}\\*
\end{longtable}

\hypertarget{taux-duxe9tagement}{%
\section{\texorpdfstring{\ul{Taux
d'étagement}}{Taux d'étagement}}\label{taux-duxe9tagement}}

Le taux d'étagement est un indiacteur qui reflète la perte de pente
naturelle due à la présence d'ouvrages transversaux. Il permet de rendre
compte de la perte de continuité longitudinale induite par des ruptures
artificielles. Le seuil maximum, fixé par le SDAGE 2022/2027, à ne pas
dépasser est de 30\%. Cependant, il est précisé que le taux d'étagement
doit être le plus bas possible pour les masses d'eau à enjeux. Ce taux
d'étagement est calculé comme suit : \[ TE = \sum_{HC}\frac{1}{pt} \]

\begin{longtable}[l]{rrr}
\toprule
État & Hauteurs cumulées & Rapport des hauteurs cumulées sur la pente naturelle (en \%)\\
\midrule
\endfirsthead
\multicolumn{3}{@{}l}{\textit{(continued)}}\\
\toprule
État & Hauteurs cumulées & Rapport des hauteurs cumulées sur la pente naturelle (en \%)\\
\midrule
\endhead

\endfoot
\bottomrule
\endlastfoot
\cellcolor{gray!15}{Existant} & \cellcolor{gray!15}{8.60} & \cellcolor{gray!15}{68.09}\\
Détruit partiellement & 4.03 & 31.91\\
\cellcolor{gray!15}{Détruit entièrement} & \cellcolor{gray!15}{0.00} & \cellcolor{gray!15}{0.00}\\*
\end{longtable}

\hypertarget{carte-du-cours-deau}{%
\section{\texorpdfstring{\ul{Carte du cours
d'eau}}{Carte du cours d'eau}}\label{carte-du-cours-deau}}

\includegraphics{test_files/figure-latex/Carte CE-1.pdf}

\end{document}
